\documentclass[12pt]{article}
\usepackage[a4paper, margin=1in]{geometry}
\usepackage{amsmath, amssymb, minted, enumitem, cmbright, hyperref}
\renewcommand{\thesection}{\Alph{section}}
\renewcommand{\thesubsection}{\thesection.\arabic{subsection}}
\renewcommand{\t}{\texttt}
\begin{document}
\section{MCQs}
The answers are \textbf{cec}...
\begin{enumerate}
\item The loop goes over every pair $(i,j)$ with $i<j$. There are $n(n+1)/2\in\Theta\left(n^2\right)$ such pairs and hence the correct answer is c). $\checkmark$
\item This code prints \t{outcome A} if there are duplicate values in the list and \t{outcome B} otherwise. In all options a)-d) there exist at least one duplicate pair, so the correct answer is e). $\checkmark$
\item The list comprehension takes $O(n)$ for $a$ and $O(n\log n)$ for $b$, and the loop is also $O(n)$ as $b$ is of constant size ($\le7$). Thus the correct answer is c). $\checkmark$
\end{enumerate}
All other questions are redacted.
\section{Simpler Questions}
\subsection{Big-O Time Complexity Analysis}
Noting the asymptotic nature of big-O notation, 1.-4. are all false. But it is stated that there is no randomized component in this unknown algorithm. Thus 5. is always true.
\subsection{A Sorting Algorithm}
This is essentially the deterministic version of quicksort.

1. is false as the elements in a Python list might not be comparable. 3. is false as it is not randomized: \t{l[0]} is always the pivot. 5. is false as it is $O\left(n^2\right)$ in that case.

Only 2. and 4. are true. (Strictly 4. could also be considered wrong as the algorithm is deterministic and `expected time complexity' is usually the term for randomized algorithms; it depends on the distribution of the input. However quicksort also enjoys average time complexity of $O(n\log n)$ so I consider it correct.)
\section{Application Questions}
\subsection{Buddy System}
\subsubsection{Check Your Understanding}
\begin{enumerate}
\item $S_1=0+1+2+3+4+5+6=21$ as each student’s buddy is the previous student.
\item $S_2=0+0+3+3+0+2+6=14$.
\end{enumerate}
\subsubsection{Subtask 1: Naïve Algorithm}
For each student $i$, scan left from $i-1$ down to $0$ and stop at the first height strictly smaller than them. It can be easily seen that this is $O\left(n^2\right)$.
\subsubsection{Solve The Full Problem}
We can use a monotonic stack to achieve $O(n)$.

The idea is as follows:
\begin{itemize}
\item Traverse $L$. For student with height $h$, remove from the stack every element $\ge h$, because such students can never be the buddy for them or indeed any later student.
\item After popping, if the stack is non-empty, its top element is the height of their buddy, so add that height to the checksum.
\item Finally, push $h$ into the stack.
\end{itemize}
\subsection{BST Merging}
\subsubsection{Subtask 1: Naïve Algorithm}
The naïve idea would be to perform inorder traversal on both BSTs, merge the two sorted lists, and create the new BST from there.

(This is essentially equivalent to IT5003 24S2 B.4.)
\subsubsection{Solve The Full Problem}
The problem states that all integer values in $A$ are strictly smaller than $B$ and also hints that we can reuse $A$ or $B$, so via the BST property we can simply attach $B$ to the rightmost node of $A$ (or equivalently, attach $A$ to the leftmost node of $B$).

Finding the rightmost node in $A$ takes $O\left(H_A\right)$ and finding the rightmost node in $B$ takes $O\left(H_B\right)$, so the algorithm is $O\left(\min(H_A+H_B)\right)$ if we choose the cheaper option.

The trick is to use an `interleaved traversal': we simultaneously move right in $A$ and move left in $B$ and stop as soon as one finishes.
\subsection{Binge Reading Competitive Programming Book}
\subsubsection{Store the Input Information}
\begin{enumerate}
\item DAG (directed acyclic graph)
\item its out-degree is $0$ (i.e., it has no outgoing edges)
\item \href{https://en.wikipedia.org/wiki/Transpose_graph}{transpose (converse/reverse)}
\item adjacency list
\item array
\end{enumerate}
\subsubsection{Check Your Understanding}
$C,D,E$; $26$ (by choosing $C$ and $E$ and reading $B,C,E$)
\subsubsection{Subtask 1: Just one “ultimate topic”}
If the chosen ultimate topic is $t$, then we can just run DFS/BFS from $t$ in $G'$ and sum the pages of all visited nodes since the reachable nodes from $t$ in $G'$ are exactly the sections that $t$ depends on. This is obviously $O(n+m)$.
\subsubsection{Subtask 2: Two specific “ultimate topics”}
Now we want the pages in the union of the prerequisites of $t_1$ and $t_2$. The idea is as follows:
\begin{itemize}
\item Run DFS/BFS from $t_1$ in $G'$ and maintain a visited array.
\item Run DFS/BFS from $t_2$ in $G'$, marking into the same visited array.
\item Sum each node only the first time it is visited.
\end{itemize}
This is still $O(n+m)$ as each node and each edge is processed at most a constant number of times.
\subsubsection{Solve The Full Problem}
Let the set of ultimate topics be $\{t_1,t_2,\dots,t_k\}$. For each ultimate topic $t_i$, we need to run DFS/BFS in the reversed graph, mark every reachable section and compute the total pages of those reachable sections. In this reachable set, for every pair of ultimate topics $(t_i,t_j)$, we need to take the union of their reachable sets, compute the total pages in that union and output the minimum.

Thus the time complexity is $O\left(k(n+m)+k^2n\right)$. As $k\in O(n)$ and $m\in O\left(n^2\right)$ in the worst case, this is equivalent to $O\left(n^3\right)$.

(Actually, the traversal part can be improved with DAG DP over bitsets, but this is out of syllabus and the asymptotic time complexity remains the same, so I omit it here but I include a Python code in the appendix.)
\newpage
\section*{Appendix: Python Code for Application Questions}
\begin{minted}{python}
from typing import List, Optional, Tuple


# Definition for a BSTvertex.
class BSTVertex:
    def __init__(self, value=0, left=None, right=None):
        self.value = value
        self.left = left
        self.right = right


def printBST(root, level=0):
    if root is not None:
        printBST(root.right, level + 1)
        print('    ' * level + str(root.value))
        printBST(root.left, level + 1)


class Solution:
    def buddySystem_naive(self, L: List[int]) -> int:
        s = 0
        for i in range(len(L)):
            for j in range(i - 1, -1, -1):
                if L[j] < L[i]:
                    s += L[j]
                    break
        return s

    def buddySystem(self, L: List[int]) -> int:
        s = 0
        stack = []
        for h in L:
            while stack and stack[-1] >= h:
                stack.pop()
            if stack:
                s += stack[-1]
            stack.append(h)
        return s

    def BSTMerging_naive(self, A: Optional[BSTVertex],
                         B: Optional[BSTVertex]) -> Optional[BSTVertex]:
        def inorder(root, arr):
            if not root:
                return
            inorder(root.left, arr)
            arr.append(root.value)
            inorder(root.right, arr)
        arrA, arrB = [], []
        inorder(A, arrA)
        inorder(B, arrB)

        def build(arr):
            if not arr:
                return None
            mid = len(arr) // 2
            root = BSTVertex(arr[mid])
            root.left = build(arr[:mid])
            root.right = build(arr[mid + 1:])
            return root
        return build(arrA + arrB)

    def BSTMerging(self, A: Optional[BSTVertex],
                   B: Optional[BSTVertex]) -> Optional[BSTVertex]:
        if not A:
            return B
        if not B:
            return A
        currA = A
        currB = B
        while currA.right and currB.left:
            currA = currA.right
            currB = currB.left
        if not currA.right:
            currA.right = B
            return A
        else:
            currB.left = A
            return B

    def bingeReading(self, n: int, m: int, pages: List[int],
                     deps: List[Tuple[int, int]]) -> int:
        g = [[] for _ in range(n)]
        o = [0] * n
        for u, v in deps:
            g[v - 1].append(u - 1)
            o[u - 1] += 1
        S = [i for i in range(n) if o[i] == 0]
        need = []
        for s in S:
            v = [0] * n
            stack = [s]
            v[s] = 1
            while stack:
                u = stack.pop()
                for p in g[u]:
                    if not v[p]:
                        v[p] = 1
                        stack.append(p)
            need.append(v)
        return min(sum(pages[x] for x in range(n) if need[i][x] or need[j][x])
                   for i in range(len(S)) for j in range(i + 1, len(S)))

    # Alternative version via bitmask DP
    def bingeReading_alt(self, n: int, m: int, pages: List[int],
                         deps: List[Tuple[int, int]]) -> int:
        g = [[] for _ in range(n)]
        o = [0] * n
        for u, v in deps:
            g[u - 1].append(v - 1)
            o[u - 1] += 1
        need = [1 << i for i in range(n)]
        for u in range(n):
            for v in g[u]:
                need[v] |= need[u]
        S = [i for i in range(n) if o[i] == 0]
        ans = 10 ** 10
        for i in range(len(S)):
            for j in range(i + 1, len(S)):
                mask = need[S[i]] | need[S[j]]
                total = sum(pages[k] for k in range(n) if (mask >> k) & 1)
                ans = min(ans, total)
        return ans


tests = [[1, 2, 3, 4, 5, 6, 7], [8, 3, 9, 5, 2, 6, 9]]
for t in tests:
    print(f"{t} -> {Solution().buddySystem_naive(t)}")
    print(f"{t} -> {Solution().buddySystem(t)}")
"""[1, 2, 3, 4, 5, 6, 7] -> 21
[8, 3, 9, 5, 2, 6, 9] -> 14"""

root1 = BSTVertex(53)
root1.left = BSTVertex(3)
root1.left.right = BSTVertex(52)
root1.left.right.left = BSTVertex(27)
root2 = BSTVertex(81)
root2.left = BSTVertex(67)
root2.left.right = BSTVertex(69)
root2.right = BSTVertex(89)
printBST(root1)
print("\n")
printBST(root2)
print("\n")
printBST(Solution().BSTMerging_naive(root1, root2))
print("\n")
printBST(Solution().BSTMerging(root1, root2))
"""
53
        52
            27
    3


    89
81
        69
    67


        89
    81
        69
67
        53
    52
        27
            3


        89
    81
            69
        67
53
        52
            27
    3
"""

tests = [[4, 2, [10, 7, 5, 2], [(1, 3), (1, 4)]],
         [5, 2, [10, 7, 4, 6, 15], [(1, 4), (2, 3)]],
         [5, 0, [8, 3, 10, 2, 7], []]]
for t in tests:
    print(f"{t} -> {Solution().bingeReading(*t)}")
    print(f"{t} -> {Solution().bingeReading_alt(*t)}")
"""[4, 2, [10, 7, 5, 2], [(1, 3), (1, 4)]] -> 17
[5, 2, [10, 7, 4, 6, 15], [(1, 4), (2, 3)]] -> 26
[5, 0, [8, 3, 10, 2, 7], []] -> 5"""
\end{minted}
\end{document}
